
\begin{titlepage}
    \centering

    {\largefont \textbf{RESUMO}}
    
    \begin{justify}
      \begin{singlespace}
        \normalsizefont{
Com o desenvolvimento da computação móvel, projetos são pensados com foco nas qualidades deste tipo de sistema. Em contraste com aplicações web, o estado offline não é um erro ou algo temporário, mas fato que o desenvolvimento móvel deve relevar. Assim o paradigma Offline-First se consolida como uma realidade; este paradigma descreve, dentre outros aspectos, a necessidade de sincronização de dados dos dispositivos que se conectam ao sistema, para manter a consistência e disponibilidade desses dados. Em cenários onde há centralização dos processos, como em um servidor web, existem estratégias específicas para a manutenção dos arquivos e estados; e o mesmo se aplica à computação descentralizada. Redes Peer-to-Peer são uma boa alternativa a um servidor devido o seu baixo custo e privacidade, atualmente ganharam mais tração com o surgimento das blockchains. No entanto, redes Peer-to-Peer apresentam seus próprios desafios de sincronização: peer churn, consistência entre as réplicas e eficiência na sincronização de dados. Desse modo, a necessidade do estudo de estratégias de sincronização em redes Peer-to-Peer em sistemas Offline-First é importante para a reflexão e o direcionamento na escolha da implementação da sincronização de dados congruente ao sistema em produção. Este estudo deve ser realizado com o uso de métricas em ambiente controlado com nós funcionando em rede com conexão descentralizada para obter dados aptos a serem analisados de modo a serem dispostos a comparação das vantagens, limitações e comportamentos dos diferentes tipos de sincronização. No mesmo sentido, o desenvolvimento de uma aplicação Offline-First real para proporcionar substâncias a cada uma dessas aplicações é indispensável para validar as discussões realizadas sobre cada estratégia de sincronização de dados.
} \\ [1cm]
    \normalsizefont \textbf{Palavras-chave: } {Offline-First, Peer-to-Peer, Sincronização de dados, CRDT, CAP.}
    \end{singlespace}
  \end{justify}

    \vfill
\end{titlepage}
